\section{Methods}

\subsection{Prospective freezing and paired construction}
Each study freezes its world generator, seed, arms, metrics, gates, tie-breaks, and honest terminal vocabulary before result-bearing execution. Episodes are paired where relevant so competing arms face the same underlying realization.

\subsection{Matched information}
Within a family, every non-oracle arm receives the identical serialized visible state. Mechanism comparisons therefore change the decision rule, not the available facts. Oracle arms are included only as upper/reference bounds and are not described as deployable baselines.

\subsection{First right of refusal}
The strongest donor-complete comparator appropriate to each construction is registered in advance. A donor tie or win closes the mechanism/policy claim. This rule is responsible for the N1-C and original-world N2-F5B negatives reported below.

\subsection{Hostile validity gates}
Worlds include a regime that should punish the shortcut under test: blind optimism, global reopening, partial chain inspection, decision-irrelevant information gain, stale carry-forward, or a typed operation that should provide no value. The expected trap must actually bite where prespecified; otherwise the world terminates invalid rather than yielding a positive mechanism claim.

\subsection{Determinism and authority}
The N4 studies use frozen seeds and deterministic decision paths. Canonical receipts encode result and scope. Arms labelled as language-model proxies refer to fixed heuristics, not to measurements of any deployed model. Authority strings explicitly limit the results to exact-synthetic mechanism isolation.